\section*{Chapitre \ref{ch:b1:electronic-effects} --- Effets électroniques et intermédiaires réactifs}

\begin{solution}{exo:b1:electronic-effects:1}
2-Méthylbutane, \ce{CH3-CH(CH3)-CH2-CH3}~: les trois \ce{CH3} sont
primaires, le \ce{CH2} secondaire, le CH tertiaire.
2-Bromo-2-méthylpropane~: les trois \ce{CH3} sont primaires, le carbone
central tertiaire. Parmi les quatre bromures \ce{C4H9Br} (1-bromobutane,
2-bromobutane, 1-bromo-2-méthylpropane, 2-bromo-2-méthylpropane), le
dernier donne le \omterm{def:b1:electronic-effects:intermediates}{carbocation} tertiaire, le plus stable.
\end{solution}

\begin{solution}{exo:b1:electronic-effects:2}
$10^{4.76 - 2.87} = 10^{1.89} = 78$; $10^{4.76 - 0.52} = 10^{4.24} =
\num{1.7e4}$.
\end{solution}

\begin{solution}{exo:b1:electronic-effects:3}
Éthanoate~: le C=O et le \ce{C-O^-} échangés entre les deux oxygènes.
4-Nitrophénolate~: le doublet de l'oxygène forme un C=O, les liaisons
doubles du cycle se déplacent, le carbone qui porte \ce{NO2} forme une
liaison C=N, et un doublet de N=O passe sur l'oxygène~: une formule avec
C=O à une extrémité du cycle, C=N à l'autre et les deux oxygènes du
groupe nitro négatifs (une formule quinoïde). Cation allyle~:
\ce{CH2=CH-CH2+} $\leftrightarrow$ \ce{+CH2-CH=CH2}.
\end{solution}

\begin{solution}{exo:b1:electronic-effects:4}
Seconde flèche~: de la liaison O--H de l'eau vers son oxygène~; produits
\ce{H2O} (issu de \ce{OH-}, désormais neutre) et \ce{OH-} (l'oxygène de
l'ancienne molécule d'eau garde le doublet, charge $-1$). Pour
l'ammoniac~: une flèche du \omterm{def:b1:lewis-resonance-vsepr:lewis}{doublet non liant} de N vers un H de
\ce{H3O+}, et une de cette liaison O--H vers l'oxygène~: N passe à
$+1$, O à 0.
\end{solution}

\begin{solution}{exo:b1:electronic-effects:5}
Éthanol (15.9) $<$ eau (14.0) $<$ phénol (9.99, délocalisation dans le
cycle) $<$
4-nitrophénol (7.16, \ce{NO2} prend la charge) $<$ acide éthanoïque
(4.76, deux oxygènes équivalents) $<$ acide trichloroéthanoïque (0.51,
trois chlores $-I$).
\end{solution}

\begin{solution}{exo:b1:electronic-effects:6}
Aniline (4.60) $<$ pyridine (5.23) $<$ ammoniac (9.25) $<$ méthylamine
(10.66). Dans l'aniline, le \omterm{def:b1:lewis-resonance-vsepr:lewis}{doublet non liant} est partagé avec le cycle
(formules avec C=N$^+$ et la charge négative en \textit{ortho} et en
\textit{para})~; protoner N ferait perdre cette délocalisation.
\end{solution}

\begin{solution}{exo:b1:electronic-effects:7}
Dans le 3-nitrophénolate, les \omterm{def:b1:lewis-resonance-vsepr:resonance}{formules mésomères} placent la charge sur
les carbones en \textit{ortho} et en \textit{para} de l'oxygène, jamais
sur le carbone qui porte \ce{NO2}~: aucune aide mésomère, d'où une
acidité plus faible que celle de l'isomère 4. Le groupe nitro attire
quand même les électrons par les liaisons $\sigma$ ($-I$), d'où une
acidité plus forte que celle du phénol.
\end{solution}

\begin{solution}{exo:b1:electronic-effects:8}
\ce{CH3+} $<$ \ce{CH3CH2+} $<$ \ce{CH2=CH-CH2+} $<$ \ce{C6H5CH2+}
$\approx$ \ce{(CH3)3C+} $<$ \ce{CH3OCH2+}~: les groupes alkyle cèdent
des électrons~; les groupes allyle et benzyle délocalisent la charge
(deux et quatre formules)~; le doublet de l'oxygène donne une formule où
chaque atome a son octet. L'ordre au milieu de la liste dépend du
milieu.
\end{solution}

\begin{solution}{exo:b1:electronic-effects:9}
\ce{OH-} (\omterm{def:b1:electronic-effects:nucleophile}{nucléophile}) attaque le carbone de \ce{CH3Br} (\omterm{def:b1:electronic-effects:nucleophile}{électrophile}),
le doublet de C--Br partant sur Br. \ce{H2O} cède un doublet à \ce{H+}
(on parle ici d'acide et de base). \ce{CN-} attaque le carbone du
carbonyle de l'éthanal, le doublet $\pi$ de C=O passant sur l'oxygène.
\ce{NH3} cède son \omterm{def:b1:lewis-resonance-vsepr:lewis}{doublet non liant} à l'orbitale vide du bore dans
\ce{BF3}.
\end{solution}

\begin{solution}{exo:b1:electronic-effects:10}
Non~: $K_a = 10^{-0.51} = 0.31$ est bien plus grand que $C$~; l'acide
est presque entièrement dissocié. $x^2/(0.010 - x) = 0.31$ donne
$x = \qty{9.7e-3}{mol/L}$, $\mathrm{pH} = 2.01$, \omterm{def:b1:predominant-reaction:dissociation}{coefficient de
dissociation} \qty{97}{\percent}.
\end{solution}

\begin{solution}{exo:b1:electronic-effects:11}
Les groupes alkyle repoussent les électrons ($+I$) vers l'oxygène de
l'alcoolate et rendent sa charge moins stable, d'autant plus qu'ils sont
nombreux. Les alcoolates plus gros sont aussi moins bien solvatés par
l'eau, qui stabilise mieux un petit oxygène chargé qu'un oxygène enfoui
parmi des groupes alkyle. Les deux effets élèvent le $\mathrm{p}K_a$.
\end{solution}

\begin{solution}{exo:b1:electronic-effects:12}
Dans un amide, le \omterm{def:b1:lewis-resonance-vsepr:lewis}{doublet non liant} de l'azote est partagé avec le C=O
(formule \ce{N+=C-O^-})~; le céder à un proton ou à un \omterm{def:b1:electronic-effects:nucleophile}{électrophile}
détruirait cette stabilisation. En milieu \omterm{def:b1:predominant-reaction:strong-acid}{acide fort}, l'amide est
protoné sur l'oxygène, qui porte la charge partielle négative.
\end{solution}

\begin{solution}{pb:b1:electronic-effects:1}
\textbf{1.} $10^{1.89} = 78$; $10^{1.52} = 33$; $10^{0.84} = 6.9$.
\textbf{2.} Non~: chaque chlore supplémentaire aide moins. La charge du
carboxylate est déjà bien étalée, et les mesures deviennent moins
précises à mesure que l'acide approche de la dissociation totale.
\textbf{3.} Chaque Cl attire la densité électronique à travers les
liaisons C--C et C--O et stabilise le carboxylate ($-I$).
\textbf{4.} L'\omterm{def:b1:electronic-effects:inductive}{effet inductif} s'éteint en deux ou trois liaisons.
\textbf{5.} Les deux acides paraissent aussi forts l'un que l'autre
(0.52 et 0.51). Dans l'eau, tous deux sont presque entièrement dissociés
aux concentrations usuelles (\cref{ch:b1:predominant-reaction}, \omterm{def:b1:predominant-reaction:levelling}{effet
nivelant})~: un $\mathrm{p}K_a$ voisin de 0 est la limite de ce que l'eau
peut distinguer, et les deux valeurs comportent des incertitudes de
quelques dixièmes.
\textbf{6.} Acide éthanoïque~: forme acide (\qty{98}{\percent})~;
chloroéthanoïque~: les deux, l'anion prédominant légèrement
($10^{3 - 2.87} = 1.3$)~; les trois autres~: anion.
\textbf{7.} \ce{CH3-C(=O)-O^-} $\leftrightarrow$ \ce{CH3-C(-O^-)=O}.
\textbf{8.} Deux liaisons C--O égales, intermédiaires entre une liaison
double et une liaison simple.
\textbf{9.} La charge est sur le seul oxygène, sans liaison $\pi$ pour la
partager.
\textbf{10.} $10^{15.9 - 4.76} = 10^{11.1}$, environ $10^{11}$.
\textbf{11.} La charge du phénolate s'étale sur trois carbones du cycle,
moins électronégatifs que l'oxygène~: un gain, plus faible que dans un
carboxylate~: $\mathrm{p}K_a$ 9.99, entre 4.76 et 15.9.
\textbf{12.} $K = 10^{6.37 - \mathrm{p}K_a}$~: acide éthanoïque
$10^{1.61} = 41$, réaction favorable (du dioxyde de carbone se dégage
avec un excès d'hydrogénocarbonate)~; phénol $10^{-3.62}$ et éthanol
$10^{-9.5}$~: pas de réaction.
\textbf{13.} 4-Nitrophénol (7.16) $>$ 2-nitrophénol (7.23) $>$
3-nitrophénol (8.36) $>$ phénol (9.99).
\textbf{14.} La formule quinoïde de l'\cref{exo:b1:electronic-effects:3},
avec la charge sur un oxygène du groupe nitro.
\textbf{15.} La charge n'atteint que les carbones en \textit{ortho} et en
\textit{para} de l'oxygène~; le carbone du groupe nitro est en
\textit{méta}.
\textbf{16.} L'effet $-I$ de \ce{NO2}.
\textbf{17.} Dans le 2-nitrophénol, le OH forme une \omterm{def:b1:intermolecular-forces-solvents:hbond}{liaison hydrogène}
avec un oxygène du groupe nitro voisin, ce qui stabilise la forme acide
et rend le proton un peu plus difficile à arracher.
\textbf{18.} Fractions anioniques $1/(1 + 10^{\mathrm{p}K_a - 7.5})$~:
4-nitro
\qty{69}{\percent}, 2-nitro \qty{65}{\percent}, 3-nitro
\qty{12}{\percent}. L'acide et l'anion du 4-nitrophénol absorbent
différemment, l'anion dans le visible~: sa couleur apparaît autour de son
$\mathrm{p}K_a$.
\textbf{19.} $x^2/(0.010 - x) = 10^{-0.52} = 0.30$~: $x = \qty{9.7e-3}{mol/L}$,
$\mathrm{pH} = 2.01$~: presque un \omterm{def:b1:predominant-reaction:strong-acid}{acide fort}.
\textbf{20.} $\frac12(4.76 + 2.00) = 3.38$.
\textbf{21.} \ce{CF3COOH + CH3COO- -> CF3COO- + CH3COOH}, $K = 10^{4.76 -
0.52} = 10^{4.24}$~: \omterm{def:b1:extent-q-and-k:quantitative}{réaction quantitative}.
\textbf{22.} $K_a(\ce{CF3COOH})/K_a(\ce{CH3COOH}) =$ \textbf{$10^{4.24} =
\num{1.7e4}$}, environ dix mille.
\end{solution}
